\documentclass{resume} % Use the custom resume.cls style
\usepackage{enumitem} % for item spacing
\usepackage{scalerel} % for inline images

\usepackage[left=0.4 in,top=0.4in,right=0.4 in,bottom=0.4in]{geometry} % Document margins
\newcommand{\tab}[1]{\hspace{.2667\textwidth}\rlap{#1}} 
\newcommand{\itab}[1]{\hspace{0em}\rlap{#1}}
\name{JOHN GIORSHEV} % Your name

\address{
\scalerel*{\includegraphics{images/email.png}}{|} \href{mailto:johngiorshev@gmail.com}{johngiorshev@gmail.com}\space\space\space\space\space
\scalerel*{\includegraphics{images/phone.png}}{|} 705-816-4454\space\space\space\space\space\space
\scalerel*{\includegraphics{images/linkedin.png}}{|} \href{https://www.linkedin.com/in/johngiorshev/}{johngiorshev}
}  % Your phone number, email, linkedin, and (optional) website

\begin{document}

\begin{rSection}{EMPLOYMENT EXPERIENCE}
    \begin{itemize}
        \item {\bf Software Intern} - SpaceRyde \hfill {Sep 2021 - Sep 2022}
        \begin{itemize}[topsep=-10pt]
            \setlength\itemsep{-0.35em}
            \item[\textbullet] {\bf Telemetry} - Created a module for prioritization, \href{http://libtins.github.io/}{packetization}, and streaming of sensors and metrics.
            \item[\textbullet] {\bf Simulator} - Created a 6DOF rigid body simulator for software and controls verification.\\
            Integrated with ROS, GeographicLib resources, threaded for concurrency, and \href{https://github.com/jbeder/yaml-cpp}{massively configurable}.
            \item[\textbullet] {\bf Hardware Abstractions} - Integrated finite state machines and QT GUIs over various components. 
        \end{itemize}
        \item {\bf Cyber Threat Intelligence Engineering Co-op} - Bell Canada Enterprises \hfill {Summer 2020}
        \begin{itemize}[topsep=-10pt]
            \setlength\itemsep{-0.35em}
            \item[\textbullet] Integrated a MISP server with an Elasticsearch database using RESTful APIs.\\
            Collaborated using Confluence and Jira. Used Docker, Docker Compose, Kafka.
            \item[\textbullet] Mattermost (slack) notification bot for faster dissemination of Indicators of Compromise.
        \end{itemize}
    \end{itemize}
\end{rSection}

\begin{rSection}{Projects}
    \begin{itemize}
        \item {\bf Key Value} - \href{https://keyvalue.duckdns.org/}{keyvalue.duckdns.org}\\
        Self-hosted full-stack web app for users to store key value pairs. RUST http backend integrated with SQLite server. TLS termination using NGINX and Let's Encrypt for secure ingress.
        Statically served JS, CSS, HTML front end. Docker Compose. 
        \item {\bf Sparse Distributed Representation} - \href{https://github.com/jagprog5/SDR/}{jagprog5/sdr}\\
        C++ sparse math library. Template metaprogramming and specializations for data structures. Integrated with CMake, CI/CD, static analyzers, fuzzing for complete code coverage and profile guided optimization.
        \item {\bf Conway App} - \href{https://github.com/jagprog5/conway_app/}{jagprog5/conway\_app}\\
        WASM based web and desktop GUI for Conway's Game of Life.
        \item {\bf Choose} - \href{https://github.com/jagprog5/choose/}{jagprog5/choose}\\
        C++ Ncurses utility for creating selection dialogs.
        Uses the PCRE2 regex engine and stl unicode libraries.
        \item {\bf Data Analytics}\\
        Implemented multi-layer perceptrons using \href{https://github.com/jagprog5/mlp}{pytorch} and in \href{https://github.com/jagprog5/CNeuralNet}{C99}, \href{https://github.com/jagprog5/hopfield}{hopfield networks}.\\
        Created a \href{https://github.com/jagprog5/CSV-Plotter}{shell wrapper} around Matplotlib for automated graphing and screenshotting of csv data. 
        \item {\bf Game Dev} - \href{https://github.com/jagprog5/GameDemo}{jagprog5/game\_demo}\\
        Java based game engine with demo. Used Photoshop, GIMP, and Audacity for asset creation.
    \end{itemize}
\end{rSection}

\begin{rSection}{Design Teams}
    \begin{itemize}
        \item {\bf QMind} - Queen's AI/ML Software Team \hfill {Sep 2019 - May 2020}\\
        Presented at \href{https://medium.com/qmind-ai/cucai2019-c9d5f848f5c}{CUCAI 2019}. Used OpenCV and Flask to create a computer vision API.
        \item {\bf Network Security Team} \hfill {Sep 2018 - May 2019}\\
        Penetration testing on virtual machines with Metasploit. Exploited apps in a sandboxed environment.
\end{itemize}
\end{rSection}

\begin{rSection}{EDUCATION}
    \begin{itemize}
        \item {\bf Computer Engineering} - Queen's University (Dean's Scholar) \hfill {Graduating May 2023}
    \end{itemize}
\end{rSection}

\end{document}
